% titel: Blatt 9, A4: Taylorentwicklung
% skript: Kap. 5.4 (Taylorentwicklung und Integration, S. 96)

% Teil I: Werkzeuge, kein Lösungsweg.

Taylor-/Potenzreihe von $f(x)=\frac1{1+x^2}$ um $x_0=0$, Herleitung der
arctan-Reihe durch gliedweise Integration und numerische Auswertung von
$\arctan(0.5)$ mit Python.

\begin{werkzeugkasten}[Potenzreihe der geometrischen Reihe]
Für $|q|<1$:
\begin{equation}
  \frac1{1+q}=1-q+q^2-q^3+\cdots=\sum_{k=0}^{\infty}(-1)^k q^k .
\end{equation}
Mit $q=x^2$ erhält man direkt
\begin{equation}
  \frac1{1+x^2}=\sum_{k=0}^{\infty}(-1)^k x^{2k},
  \qquad|x|<1 .
\end{equation}
\textit{Skript:} geometrische Reihe, Kap. 2.1; Potenzreihen, Kap. 5.4.
\end{werkzeugkasten}

\begin{werkzeugkasten}[Gliedweise Integration]
In ihrem Konvergenzintervall darf eine Potenzreihe gliedweise integriert
werden (Satz 5.4.2):
\begin{equation}
  \int_0^x \sum_{k=0}^{\infty}a_k t^k\,\dd t
  = \sum_{k=0}^{\infty}a_k\frac{x^{k+1}}{k+1}.
\end{equation}
Wendet man das auf die Reihe für $\frac1{1+t^2}$ an, erhält man eine Reihe für
$\arctan$, weil $\arctan'(t)=\frac1{1+t^2}$ und $\arctan(0)=0$.

\textit{Skript:} Satz 5.4.2, Beispiel 5.4.5.
\end{werkzeugkasten}

\begin{werkzeugkasten}[Numerik mit Python]
Eine Potenzreihe mit endlichen Summen approximieren: Reihe bis $N$ aufsummieren,
mit einem Summationsschritt der als Bruch oder via $\mathrm{math}$ exakt bleibt.
Dezimalausgabe z. B. mit `print(...)`. Alternierende Reihen konvergieren oft
schnell und der Fehler ist durch das erste weggelassene Glied abgeschätzt.

\textit{Skript:} Anhang A.2 (einige Python-Befehle), S. 139.
\end{werkzeugkasten}

\begin{achtung}[Darauf achten]
\begin{itemize}
  \item Die Reihe für $\frac1{1+x^2}$ hat nur gerade Exponenten; der
        Koeffizient $a_{2k}=(-1)^k$, alle ungeraden Koeffizienten sind $0$.
  \item Beim Integrieren jedes Glied $x^{2k}\to\frac{x^{2k+1}}{2k+1}$ — den
        neuen Nenner $2k+1$ nicht vergessen.
  \item Die Reihen konvergieren absolut für $|x|<1$; für $\arctan(0.5)$ ist
        die Konvergenz gut.
  \item In Python Potenzen sauber schreiben (z. B. `(-1)**k * x**(2*k+1) /
        (2*k+1)`) und genügend Glieder summieren.
\end{itemize}
\end{achtung}
